\documentclass[a4paper, 12pt, oneside, headspline]{scrartcl}
% !TeX root = ../main.tex
%Variablen welche innerhalb der gesamten Arbeit zur Verfügung stehen sollen
\newcommand{\titleDocument}{Bachelorarbeit}
\newcommand{\subjectDocument}{Test-driven development mit Python}
\newcommand{\institution}{Hochschule für angewandte Wissenschaften Augsburg}
\newcommand{\insSection}{Fakultät für Informatik}
\newcommand{\name}{Maximilian Konter}
\newcommand{\email}{maximilian.konter@hs-augsburg.de}


% Farbige links
\usepackage{color}

% Umlaute unter UTF8 nutzen
\usepackage[utf8]{inputenc}

% deutsche Silbentrennung
\usepackage[ngerman]{babel}

% Grafiken aus PNG Dateien einbinden
\usepackage{graphicx}

% bricht lange URLs "schoend" um
\usepackage[hyphens,obeyspaces,spaces]{url}

% mehrseitige Tabellen ermöglichen
\usepackage{longtable}

% für Tabellen
\usepackage{array}

% Schaltet den zusätzlichen Zwischenraum ab, den LaTeX normalerweise nach einem Satzzeichen einfügt.
\frenchspacing

% Paket für Zeilenabstand
\usepackage{setspace}

% für Bildbezeichner
\usepackage{capt-of}

% für Stichwortverzeichnis
\usepackage{makeidx}

% Festlegung Art der Zitierung - Havardmethode: Abkuerzung Autor + Jahr
\bibliographystyle{alphadin}

% für Listings
\usepackage{listings}
\lstset{numbers=left, numberstyle=\tiny, numbersep=5pt, keywordstyle=\color{black}\bfseries, stringstyle=\ttfamily,showstringspaces=false,basicstyle=\footnotesize,captionpos=b}
\lstset{language=python}

% Indexerstellung
\makeindex

% Abkürzungsverzeichnis
\usepackage[german]{nomencl}
\let\abbrev\nomenclature

% Abkürzungsverzeichnis LiveTex Version
\renewcommand{\nomname}{Abkürzungsverzeichnis}
\setlength{\nomlabelwidth}{.25\hsize}
\renewcommand{\nomlabel}[1]{#1 \dotfill}
\setlength{\nomitemsep}{-\parsep}
\makenomenclature

% Disable single lines at the start of a paragraph (Schusterjungen)
\clubpenalty = 10000
% Disable single lines at the end of a paragraph (Hurenkinder)
\widowpenalty = 10000
\displaywidowpenalty = 10000

% Hyperlinks im Text
\usepackage[bookmarksnumbered,pdftitle={\titleDocument},hyperfootnotes=false]{hyperref}
\hypersetup{colorlinks, citecolor=red, linkcolor=blue, urlcolor=black}

% neue Kopfzeilen mit fancypaket
\usepackage{fancyhdr} %Paket laden
\pagestyle{fancy} %eigener Seitenstil
\fancyhf{} %alle Kopf- und Fußzeilenfelder bereinigen
\fancyhead[L]{\nouppercase{\leftmark}} %Kopfzeile links
\fancyhead[C]{} %zentrierte Kopfzeile
\fancyhead[R]{\thepage} %Kopfzeile rechts
\renewcommand{\headrulewidth}{0.4pt} %obere Trennlinie
\fancyfoot[C]{\thepage} %Seitennummer
\renewcommand{\footrulewidth}{0.4pt} %untere Trennlinie

\begin{document}
	% Leere Seite am Anfang
	\newpage
	\thispagestyle{empty} % erzeugt Seite ohne Kopf- / Fusszeile
	\section*{ }
	
	% Titelseite
	% !TeX root = ../bachelor.tex
\thispagestyle{empty}

\begin{figure}[t]
 \centering
 \includegraphics[width=0.3\textwidth]{abb/logoHSA}
\end{figure}


\begin{verbatim}


\end{verbatim}

\begin{center}
\Large{\institution}
\end{center}


\begin{center}
\Large{\insSection}
\end{center}
\begin{verbatim}




\end{verbatim}
\begin{center}
\doublespacing
\textbf{\LARGE{\titleDocument}}
\singlespacing
\begin{verbatim}

\end{verbatim}
\textbf{{~\subjectDocument~}}
\end{center}
\begin{verbatim}

\end{verbatim}
\begin{center}

\end{center}
\begin{verbatim}

\end{verbatim}
\begin{center}
\textbf{zur Erlangung des akademischen Grades \\ Bachelor of Science}
\end{center}
\begin{verbatim}





\end{verbatim}

\begin{footnotesize}
\begin{flushleft}
\begin{tabular}{llll}
\textbf{Thema:} & & \subjectDocument & \\
& & \\
\textbf{Autor:} & & \name & \\
& & \email & \\
& & MatNr. 951004 & \\
& & \\
\textbf{Version vom:} & & \today &\\
& & \\
\textbf{1. Betreuer:} & & Dipl.-Inf. (FH), Dipl.-De Erich Seifert, MA &\\
\textbf{2. BetreuerIn:} & & Prof. Dr.-Ing. Alexandra Teynor &\\
\end{tabular}
\end{flushleft}
\end{footnotesize}

	% 1.5 facher Zeilenabstand
	\onehalfspacing
	
	% Einleitung / Abstract
	% !TeX root = main.tex
	
	% einfacher Zeilenabstand
	\singlespacing
	
	% Inhaltsverzeichnis anzeigen
	\newpage
	\tableofcontents
	
	% das Abbildungsverzeichnis
	\newpage
	% Abbildungsverzeichnis soll im Inhaltsverzeichnis auftauchen
	\addcontentsline{toc}{section}{Abbildungsverzeichnis}
	% Abbildungsverzeichnis endgueltig anzeigen
	\listoffigures
	
	% das Tabellenverzeichnis
	\newpage
	% Abbildungsverzeichnis soll im Inhaltsverzeichnis auftauchen
	\addcontentsline{toc}{section}{Tabellenverzeichnis}
	% \fancyhead[L]{Abbildungsverzeichnis / Abkürzungsverzeichnis} %Kopfzeile links
	% Abbildungsverzeichnis endgueltig anzeigen
	\listoftables
	
	% das Listingverzeichnis
	\newpage
	% Listingverzeichnis soll im Inhaltsverzeichnis auftauchen
	\addcontentsline{toc}{section}{Listingverzeichnis}
	\fancyhead[L]{Listingverzeichnis} %Kopfzeile links
	\renewcommand{\lstlistlistingname}{Listingverzeichnis}
	\lstlistoflistings
	
	% das Abkürzungsverzeichnis
	\newpage
	% Abkürzungsverzeichnis soll im Inhaltsverzeichnis auftauchen
	\addcontentsline{toc}{section}{Abkürzungsverzeichnis}
	% das Abkürzungsverzeichnis entgültige Ausgeben
	\fancyhead[L]{Abkürzungsverzeichnis} %Kopfzeile links
	% !TeX root = ../../main.tex
\nomenclature{GPL}{GNU General Public License}
\nomenclature{GNU}{GNU is not Unix}
\nomenclature{LGPL}{GNU Lesser General Public License}
\nomenclature{TDD}{Test-driven development}
\nomenclature{GUI}{Graphical User Interface}
\nomenclature{STDLIB}{Python Standard Bibliothek}
\nomenclature{VCS}{\gls{vcs}}
	\printnomenclature
	
	%%%%% Beginn Dokument %%%%%
	\newpage
	\fancyhead[L]{\nouppercase{\leftmark}} %Kopfzeile links
	
	% 1,5 facher Zeilenabstand
	\onehalfspacing
	
	% einzelne Kapitel
	% !TeX root = ../main.tex
\section{Einleitung}\label{einleitung}

Das ist meine Einleitung
	
	\onecolumn
	
	% einfacher Zeilenabstand
	\singlespacing
	
	% Literaturliste soll im Inhaltsverzeichnis auftauchen
	\newpage
	\addcontentsline{toc}{section}{Literaturverzeichnis}
	
	% Literaturverzeichnis anzeigen
	\renewcommand\refname{Literaturverzeichnis}
	\bibliography{Hauptdatei}
	
	% 1,5 facher Zeilenabstand
	\onehalfspacing

	\newpage
	
	% Eidesstattliche Erklärung
	\addcontentsline{toc}{section}{Eidesstattliche Erklärung}
	% !TeX root = bachelor.tex
\newpage

\thispagestyle{empty}
\addcontentsline{toc}{section}{Eidesstattliche Erklärung}

\begin{verbatim}

\end{verbatim}

\begin{LARGE}
\begin{center}
	\textbf{Eidesstattliche Erklärung zur Abschlussarbeit}
\end{center}
\end{LARGE}
\begin{verbatim}


\end{verbatim}
Hiermit versichere ich, die eingereichte Abschlussarbeit selbständig verfasst und
keine andere als die von mir angegebenen Quellen und Hilfsmittel benutzt zu haben.
Wörtlich oder inhaltlich verwendete Quellen wurden entsprechend den anerkannten
Regeln wissenschaftlichen Arbeitens zitiert. Ich erkläre weiterhin, dass die
vorliegende Arbeit noch nicht anderweitig als Abschlussarbeit eingereicht wurde.
\\
\\
Das Merkblatt zum Täuschungsverbot im Prüfungsverfahren der Hochschule
Augsburg habe ich gelesen und zur Kenntnis genommen. Ich versichere, dass die
von mir abgegebene Arbeit keinerlei Plagiate, Texte oder Bilder umfasst, die durch
von mir beauftragte Dritte erstellt wurden.
\begin{verbatim}



\end{verbatim}


\begin{displaymath}
% use packages: array
\begin{array}{ll}
Unterschrift:~~~~~~~~~~~~~~~~~~~~~~~~~~~~~~~~~~~~~~~~~~
& Ort, Datum:~~~~~~~~~~~~~~~~~~~~~~~~~~~~~~~~~~~~~~~~~~
\end{array}
\end{displaymath}
	
	% leere Abschlussseite
	\newpage
	\thispagestyle{empty} % erzeugt Seite ohne Kopf- / Fusszeile
\end{document}